\section{总结与延伸}

本讲从 PPO 的完整系统出发，解释 GRPO 如何去掉 value model，再审计其 baseline、standard-deviation 和 length biases。DeepSeek R1、Kimi K1.5 与 Qwen3 说明，成功的 RLVR recipe 需要算法、数据、curriculum、distillation、general alignment 和 systems infrastructure 协同。

七个核心结论是：

\begin{enumerate}
\item Verifiable reward 降低了标签成本和 reward ambiguity，但不会自动保证任务定义正确。
\item PPO 的复杂度主要来自 on-policy system，而不只是一条 clipped loss。
\item GRPO 简化 value estimation，却引入 finite-group、variance 和 length weighting 问题。
\item R1-Zero 的现象应做 controlled analysis，不能仅凭单条 ``aha'' trajectory 讲 emergence 故事。
\item Distillation 与 RLVR 互补：前者便宜迁移，后者主动探索。
\item Curriculum 应随 policy 能力动态变化；静态 hard set 会迅速过时。
\item RLVR 的关键系统指标是 verified useful trajectories，而不是裸吞吐。
\end{enumerate}

下一讲转向 multimodality：如果 Transformer 只理解 tokens，图像、视频和音频怎样被编码、注入、生成，并与语言模型共同训练。

\subsection{拓展阅读}

建议按“算法定义→objective 审计→完整模型 recipe→系统与环境”顺序阅读。每篇材料都应记录 policy 初始化、prompt distribution、group size、reward/verifier、长度聚合、rollout tokens、reference/KL、独立评测与失败尝试；否则不同论文中的 ``GRPO'' 或 ``RLVR'' 很可能并不是同一个实验对象。

\begin{longtable}{L{0.22\textwidth}L{0.31\textwidth}L{0.37\textwidth}}
\toprule
材料 & 为什么值得读 & 带着什么问题读 \\
\midrule
Schulman et al., \emph{Proximal Policy Optimization Algorithms} & PPO clipped surrogate 的原始定义与 trust-region 动机 & Clipping 约束的是 sample ratio 还是全局 policy distance？ \\
AlpacaFarm PPO implementation & 把 rollout、reward shaping、value/GAE 与 minibatch update 连成可运行代码 & 哪些稳定化选择不在简化公式中，张量 mask 怎样验证？ \\
DeepSeekMath / GRPO & Group-relative advantage 与早期 reasoning RL 结果的来源 & 组均值、标准差和 token averaging 分别重写了什么权重？ \\
Liu et al., \emph{Understanding R1-Zero-Like Training}（Dr. GRPO） & 系统分析 GRPO bias、长度增长与 leave-one-out 修正 & ``Aha moment'' 的哪些现象能由 base prior 或 objective bias 解释？ \\
DeepSeek R1 report & R1-Zero、cold-start SFT、reasoning RL、general alignment 与 distillation 的组合案例 & 最终分数能证明什么，哪些阶段仍缺少可归因消融？ \\
Kimi K1.5 report & 公开 difficulty filtering、length reward、curriculum 与 RL infrastructure 细节 & 如何用 rolling success rate 维持有效学习信号，并计算 verified success cost？ \\
Qwen3 / Qwen3-Coder reports & 展示低数据 RL、thinking-mode fusion、stage composition 与 agent environments & Unique prompts、rollout compute、mid-training data 和 environment verifier 如何共同贡献？ \\
\bottomrule
\caption{Lecture 16 的注释式拓展阅读路线。}
\end{longtable}

\begin{knowledgebox}{建议的复现实验}
在一个小型数学模型上固定 prompt set 与 rollout budget，对比 vanilla GRPO、leave-one-out REINFORCE、去标准差版本和不同 token-length aggregation；同时报告 answer accuracy、group success histogram、平均长度、每题梯度范数与 wall-clock verified successes。随后加入动态 curriculum，观察收益究竟来自 objective 还是题目重新加权。
\end{knowledgebox}

\end{document}
